% Continuation of 26 November 2026 at e2968b2898b5. Current proofs and primary
% passages are recorded in the prefix and curved-filling assessment in parent reviews.
% Drafted for Edward B. Baker III with OpenAI Codex assistance, 26 November 2026.
% Author-selected lead GPT-6.2 Astra/ultra; independent serving metadata unavailable.
% Scientific inputs: ece209cc09aa31629a8b496ae9d01957e71c70ac, including both
% child note sets, the parent assessment, and the 17 pins in ../references/README.md.
% Internal master and companion citations denote working manuscripts.
\documentclass[11pt,letterpaper]{article}
\usepackage[margin=1in]{geometry}
\usepackage[T1]{fontenc}
\usepackage{lmodern,microtype,amsmath,amssymb,amsthm,mathtools}
\usepackage[hidelinks]{hyperref}
\setlength{\emergencystretch}{2em}
\numberwithin{equation}{section}
\newtheorem{proposition}{Proposition}[section]
\newtheorem{theorem}[proposition]{Theorem}
\newtheorem{corollary}[proposition]{Corollary}
\DeclareMathOperator{\Ric}{Ric}
\newcommand{\Sph}{\mathbb S}
\newcommand{\Z}{\mathbb Z}
\newcommand{\R}{\mathbb R}
\newcommand{\HS}{\mathrm{HS}}
\title{Marked collars, handle reduction and\\smooth four-sphere recognition}
\author{Edward B. Baker III}
\date{Working draft, 26 November 2026}
\begin{document}
\maketitle
\begin{abstract}
A regular level of a handle decomposition records how to reassemble a
smooth four-manifold. A sphere selected by curvature records what a metric
comparison could save. We ask when one embedded hypersurface carries
both kinds of information. For a mixed presentation with two middle
handles, positive scalar curvature on either of two specified intermediate
three-manifolds recognizes the closed homotopy sphere. The two orders
have different recognizing boundaries. For a positive four-dimensional
conformal class of subcritical chiral Weyl energy, a conformally round,
umbilic three-sphere gives a complementary conclusion. Each capped
side has positive-class energy infimum at most its original mass, in
either chirality. On a hypothetical nonstandard sphere the side which
carries its smooth type must therefore carry at least its fixed-type
infimum. This excludes such a face on each sufficiently late charged
disk selected by the master manuscript. We use these results to specify
a penultimate-cut criterion for a net cancellation at any upper rank,
a curved disk-filling test, and the geometric selection still
required to carry it out on arbitrary smooth homotopy four-spheres.
\end{abstract}

\section{The problem on the original smooth manifold}

Let \(M\) be a closed, connected, oriented smooth manifold homotopy
equivalent to \(\Sph^4\). The target is a diffeomorphism \(M\cong\Sph^4\).
Every small coordinate disk has a smoothly embedded three-sphere as
its boundary. Its exterior can retain the entire unknown smooth type.
A cut useful for recognition must either isolate a
specified handle interaction, or discard enough curvature to improve a
metric on an identified smooth type.

There are accordingly two ways to use a cut. In the smooth argument,
carry out a relative handle move and attach the original remaining
handles through the induced boundary map. Count the middle handles
after the whole move. In the metric argument, choose a metric on the
given manifold, alter it through a specified path, cap the separated
pieces, and identify the smooth type of the piece on which the curvature
comparison is used. The first argument can work at a nonspherical
level. The round cap in the second requires a spherical level with
control in its actual marking. The two arguments can also be used
independently.

The deformation paper~\cite{Baker2026CollarDeformation} starts from an
embedded sphere and estimates the cost of producing a round product.
The surgery paper~\cite{Baker2026HandleCuts} starts from a marked
handle presentation and identifies the levels at which a recognition
theorem can apply. Here we develop the implication between their
conclusions and the master program~\cite{Baker2026ChiralMaster}. Two
distinct existence questions recur: which level can be selected, and
which metrics or estimates exist there?

\subsection{What is marked}

A spherical comparison uses an embedding
\[
 \Phi:\Sph^3\times(-w_L,w_R)\longrightarrow M
\]
and the inverse on its image. Its middle sphere separates \(M\) into
two complete regions \(P_L,P_R\). A normal coordinate points from left
to right. If the cut copies are parametrized by \(a_L,a_R\), the
angular identification from left to right is \(a_Ra_L^{-1}\). The
reverse identification is \(a_La_R^{-1}\). The two outward normals
have opposite signs. A change of coordinates, including one used to
make the face round, must be composed with these maps on their actual
domains. In a handle diagram the data also include each embedding
\(S^1\times D^2\to\partial X\) for an unspent two-handle and the
attaching map of the upper body. An isotopy transports the solid tori,
and hence their framings. A bound measured in a new diagram has metric
content only after identifying the metric there.

\section{Two tests of an intermediate handle level}\label{sec:handles}

Cancel the lower and upper trees at the extrema of a Morse
presentation by the marked end reductions of~\cite{Baker2026MarkedHeights}.
Write \(\ell,r\) for the remaining numbers of one- and three-handles,
and \(n\) for the number of two-handles. There is one minimum and one
maximum, so \(\chi(M)=2\) gives \(\ell+r=n\). For \(n=0\) the two end
balls are joined by a boundary diffeomorphism that extends over a ball.
For \(n=1\), after reversal the presentation is a knot trace in a
four-ball with outgoing boundary \(S^1\times S^2\). First homology fixes
the framing at zero, and Gabai's Property R makes the knot an unknot.
The diffeomorphism of the upper boundary extends over the upper
one-handlebody. Thus these two counts already recognize \(M\)
\cite{Gabai1987FoliationsIII,LaudenbachPoenaru1972Handlebodies,Baker2026MarkedHeights}.

\paragraph{The first cancellation.}
Let a framed two-handle meet a one-handle belt sphere once transversely.
Smooth adjacent-handle cancellation gives a diffeomorphism on the
trace, with an outgoing change of boundary coordinates
\cite{Milnor1965HCobordism}. In the split control add a zero-framed
unknot as a residual tube. The canceling component first gives a ball
and carries the residual unknot into its boundary. The residual
component first gives \(\#_2(S^1\times S^2)\) as the intermediate
boundary. Sliding the residual component over a framed parallel of
the first changes its geometric position and its framing together.
In the specified slid control both components pass geometrically once
through the belt, so either first attachment gives a spherical cut.
If \(F\) is the resulting new-to-old trace diffeomorphism and
\(\psi:\partial X_{\rm old}\to\partial Q\) is the original
identification with the upper body, use \(\psi\circ\partial F\) to
close the new trace. Then \(F\cup\mathrm{id}_Q\) identifies the
closed manifolds. This also determines the inverse transport for the
reversed presentation. The two orders in this control are established
by the framed embeddings~\cite{Baker2026HandleCuts}.

There is a metric realization of the first spherical cut in this
particular control. Use the two actual ball parametrizations to
identify the sides with hemispheres, extending each discrepancy in
the equatorial map over its ball and using a collar before gluing
\cite{Hatcher1983Smale,Palais1960Extending}. In the pulled-back round
metric the collar reads \(dq^2+\cos^2 q\,\gamma\), with \(\gamma\)
the unit round three-metric. Multiplication by
\(\exp(2s\beta(q)\log\sec q)\), with \(\beta=1\) near zero and zero
near both outer faces, gives a product near the equator at \(s=1\)
in arclength \(dt=\sec q\,dq\). Every metric in the path has zero
Weyl tensor. Both entire sides and their joined peak therefore have
zero energy in both signs. The same diffeomorphism that carries the
framed tubes carries this metric. For a general diagram the metric is
a separate input.

\subsection{Two middle handles and a common belt}

Now let \((\ell,n,r)=(1,2,1)\). The two winding numbers generate
\(\Z\), since the one-handle generator must die. Genuine framed
slides and sign choices implement the Euclidean algorithm, giving
a row \((1,0)\). Denote by \(Y_a\) the boundary after the winding-one
component is attached first and by \(Y_b\) the boundary after the
winding-zero component is attached first. Call \(Y\) \emph{dominated}
here if a closed oriented three-manifold admitting positive scalar
curvature maps to \(Y\) with nonzero degree.

\begin{proposition}[The two orders, \cite{Baker2026HandleCuts}]
The following tests apply to the specified intermediate boundaries:
\[
  Y_a\text{ dominated}\ \Longleftrightarrow\ Y_a\cong S^3,
  \qquad
  Y_b\text{ dominated}\ \Longleftrightarrow\
       Y_b\cong\#_2(S^1\times S^2).
\]
Either test recognizes the original \(M\).
\end{proposition}
\begin{proof}
The first lower body has a single one/two pair with unit cellular
differential. Its attaching circle kills the generator of its
fundamental group, so the body is contractible. Its boundary is an
integral homology sphere. Chodosh, Maximo and Mukherjee show that a
dominated integral homology sphere is \(S^3\) or a connected sum of
signed Poincar\'e spheres. Every nonempty such sum is excluded as the
boundary of a smooth compact contractible four-manifold
\cite[Corollary 2.3 and Proposition 4.2]{ChodoshMaximoMukherjee2024Complete}.

For the other order, use the ordinary boundary surgery diagram obtained
by giving the dotted unknot framing zero. In the order consisting of
that unknot followed by the winding-one and winding-zero components,
the full linking matrix is
\[
 \begin{pmatrix}0&1&0\\1&u&v\\0&v&c\end{pmatrix}.
\]
Full surgery has boundary \(S^1\times S^2\) from the upper body.
The matrix therefore has zero determinant, so \(c=0\). The two components defining \(Y_b\)
now have zero linking matrix. Their ordinary two-handle trace \(V\)
is simply connected. Its pairing vanishes and \(H_2(V)=\Z^2\).
The exact sequence of the pair and duality give an isomorphism
\(H_2(Y_b)\to H_2(V)\), and \(H_1(Y_b)=\Z^2\). Under domination, the
three-manifold decomposition~\cite[Proposition 2.2]{ChodoshMaximoMukherjee2024Complete}
therefore has two \(S^1\times S^2\) summands. The spherical factors
have trivial first homology. Attach three-handles on the two cross
sections. They kill a basis of \(H_2(V)\) with no kernel, and preserve
simple connectivity. The remaining boundary is a smooth integral
homology sphere bounding a contractible manifold. Its spherical
summands and their connected sum admit positive scalar curvature
\cite[Theorem A]{GromovLawson1980Surgery}. The same homology-sphere
classification and signed-sum obstruction remove the remaining
factors. Conversely \(S^3\) and \(S^1\times S^2\) have positive
round and product metrics, respectively, and connected sum gives a
positive scalar curvature metric on \(\#_2(S^1\times S^2)\).

The trace \(V\) served to identify its boundary. Recognition of \(M\)
uses the mixed theorem of~\cite{Baker2026MarkedHeights} on the original
two pieces with their upper gluing. It applies to either of the two
displayed boundary types.
\end{proof}

This has a four-dimensional geometric test even when the cut is
nonspherical. For a compact physical product band about the chosen
level let \(d\) be the intrinsic distance between the band ends and
\(\rho=\min R_g>0\). If the level admitted no positive scalar
curvature metric, R\"ade's strict width theorem applied to
\((\rho/12)g\) in dimension four would give
\(\sqrt{\rho/12}\,d<\pi/2\). Thus \(\rho d^2\ge3\pi^2\) certifies the
corresponding recognition, including equality
\cite[Corollary 2.25]{Raede2022Bands}. The companion paper derives
tests using both end mean curvatures, the lapse of a measured slice,
or stability of a separator of nonzero projection degree
\cite{Baker2026HandleCuts}. If neither intermediate cut is accessible,
enclosing the whole lower cluster in this presentation ends at
\(S^1\times S^2\), since that is the boundary of the actual upper
one-handlebody. A spherical cap comparison there requires a different
cut.

\paragraph{The next count.}
With one lower one-handle and three middle handles, a dominated
winding-one first cut again gives \(S^3\). In the dual description the
one-component surgery to \(S^1\times S^2\) identifies the unknot by
Property R. The original attaching circle is its surgery core, up to
a change of the one-handlebody marking that extends. It cancels,
carrying both remaining framed tubes. This strictly lowers \(n\)
from three to two on the same smooth manifold. The returned diagram is
a two-component framed link in \(S^3\). Its actual upper body has two
one-handles when read in reverse, so full surgery has boundary
\(\#_2(S^1\times S^2)\). Its free rank-two first homology forces the
entire two-by-two linking matrix to vanish. If the next one-component
surgery cut is also dominated, the zero-pairing comparison just used,
in rank one, gives \(S^1\times S^2\). Property R identifies that link
component as an unknot. The two-component unlink-slide theorem in
this case and extension over the upper handlebody then recognize \(M\)
\cite[Proposition 3.2]{GompfScharlemannThompson2010}%
\cite{LaudenbachPoenaru1972Handlebodies,Baker2026HandleCuts}.
The first cancellation permits arbitrarily many moving residual tubes.
The second conclusion uses the specific two-component successor.


\subsection{A reduction at the last middle handle}
Let \(Y_r=\#_r(S^1\times S^2)\), with \(Y_0=S^3\). Suppose the
upper body has rank \(q\ge1\) and the sublevel \(W\) through all but
the last middle handle is simply connected. If its actual boundary is
dominated with nonzero degree by a closed oriented PSC three-manifold,
then the presentation on \(M\) admits a net 2/3 cancellation.
The number of earlier attachments may be arbitrary
\cite{Baker2026HandleCuts}.

Here is why the boundary suffices. The intersection form of \(W\)
vanishes in the homotopy sphere, and its \(H_2\) has rank \(q-1\).
The factor spheres in the PSC boundary decomposition give an
integral basis. Compressing them in an auxiliary filling leaves
a smooth contractible manifold bounding the remaining spherical
homology sum, so that sum is trivial
\cite[Proposition 4.2]{ChodoshMaximoMukherjee2024Complete}.
Thus the actual final attachment goes from \(Y_{q-1}\) to \(Y_q\).
Its trace group has free presentations from both ends. The upper
one has a single relator with primitive abelianization and gives
zero second group homology. Applied at the lower end, Hopf's
formula and residual nilpotence of the free group force the attaching
knot to be null-homotopic. After removing common product summands
of its exterior, Ni's zero-surgery theorem puts the knot on a
zero-framed unknot in a ball \cite{Ni2023Null,Baker2026HandleCuts}.
The dual core then meets one upper belt sphere once, after a change
of boundary marking extending over that upper handlebody. The inverse
cancellation restores the original gluing and all earlier tubes.

In the normalized one-lower, three-middle presentation \((a,b,c)\),
with \(a\) of winding one, the sublevel through \(a,b\) is simply
connected. A certificate there returns the \emph{mixed} successor
\((a,b)\) even if the Mazur-type first cut fails. A certificate after
\(a\) removes the only lower 1-handle and returns the \emph{one-sided} link of the transported
\(b,c\). The stabilized Mazur double realizes the former alternative
with its long-knot first cut nonspherical. In the unstabilized double
a closed PSC metric constructed on the actual meridian-first cut has
a degree-one stable slice and an arbitrarily prescribed product
length. Its full Weyl price grows linearly with length
\cite{Baker2026HandleCuts}. These examples start with the specified
marked presentation.

There is also a route which does not use a winding-one first cut.
With one lower 1-handle, normalize the row to \((1,0,\ldots,0)\), put all
\(m-1\) zero-winding components first, and the winding-one component last. If each successive pure prefix is PSC-dominated, the boundary
chain goes from \(Y_1\) through \(Y_m\). Each step is a split
zero-framed attachment by the same single-knot result. The last
dual attachment from \(Y_{m-1}\) to \(Y_m\) is split as well.
Interior surgery on its sphere gives an upper \(m\)-handlebody.
The identified lower piece closes with it to \(S^1\times S^3\);
the inverse surgery circle generates the fundamental group
and either framing recovers \(S^4\) \cite{Baker2026HandleCuts}.
This recognizes the whole chain. A single penultimate certificate
with a simply connected lower sublevel already gives strict descent
on a least-count presentation.

\section{The cost of making a spherical cut}\label{sec:cost}

Let \(W^\pm\) be the two real Weyl operators on self-dual and
anti-self-dual two-forms in the chosen orientation. Use the curvature
convention in which the unit sphere has identity curvature operator.
For \(\sigma\in\{+,-\}\) write
\[
 E_\sigma(g;U)=\int_U |W^\sigma_g|_{\HS}^2\,dV_g,\qquad
 T=4\pi^2,\qquad I^+(M)=\inf_{Y([g])>0}E_+(g;M).
\]
Here \(Y([g])\) is the Yamabe constant of the conformal class, and the
empty infimum is \(+\infty\). Both measures are conformally invariant:
the operator weight under \(e^{2u}g\) is \(e^{-2u}\) and the volume
weight is \(e^{4u}\). Signature gives \(E_+-E_-=12\pi^2\tau(M)=0\)
on the closed manifold~\cite{Hirzebruch1971Signature}. We write \(e\)
for that common total and \(m_j^\sigma=E_\sigma(g;P_j)\) separately
on each side. A reversal of the ambient orientation exchanges the two
chiralities. Reversing just the chosen normal interchanges the sides.

\paragraph{The face and its two widths.}
Suppose the induced metric on a smooth embedded \(S^3\) is conformal
to the round metric via a specified angular map, and its second
fundamental form \(K\) is pure trace. With \(H=\operatorname{tr}K\)
and \(K=\tfrac12\mathcal L_\nu h\) in normal coordinates, the law
\(\widehat K=e^u(K+\nu(u)h)\) prescribes just two jets of \(u\).
Take \(u\) on the sphere to make its metric round and prescribe
\(\nu(u)=-H/3\), then extend these jets in a physical collar.
In Gaussian coordinates of \(\widehat g=e^{2u}g\) we have
\(\widehat g=dt^2+h_t,\ h_0=\gamma,\ h'_0=0\).
The inverse chart gives the normal foot and signed distance. Choose
each \(w_j\) inside the available region on its own side.

Let \(0\le\chi_j\le1\) be a fixed smooth cutoff, one near \(r=0\) and zero
near \(r=1\). On side \(j\) use \(r=|t|/w_j\) and interpolate by
\[
 h_{t,s}=\gamma+(1-s\chi_j(r))(h_t-\gamma),\qquad 0\le s\le1.
\]
The two formulas agree as germs across the middle. They are convex
combinations of positive tensors and recover the incoming tensor near
each outer face. Restoring the original conformal factor gives a
path from the original metric. A localized conformal adjustment
gives an exact product germ at the endpoint. For \(h_t-\gamma=O(t^2)\)
with its differentiated Taylor bounds, the two normal derivatives
falling on the cutoff cost \(O(1)\), so every half-strip has energy
\(O_g(w_j)\) uniformly in \(s\).

The companion calculation gives the signed leading term. Write
\(A=h''_0/2\) and let \(A^\circ\) denote its trace-free part. The
trace-free normal curvature block to leading order is
\(-(1-sB_j)A^\circ\), where \(B_j=(r^2\chi_j)''/2\). The tangential
block is round and the mixed block is \(O(w_j)\). Projection onto
either Weyl block contributes a factor of \(1/2\) in the leading
operator. Since \(\int_0^1 B_j\,dr=0\), the linear term integrates away.
The resulting full-side increment is
\begin{equation}\label{eq:jet}
 \Delta_j^\sigma(s)
 = w_j s^2\,\frac{\|A^\circ\|_{L^2(\gamma)}^2}{16}
       \int_0^1 r^4|\chi_j''|^2\,dr + O_{C^2_s,g}(w_j^2).
\end{equation}
The remainder controls two derivatives in \(s\). When
\(A^\circ\not\equiv0\) the increments are strictly convex and the
positive endpoint increment is the maximum for sufficiently small
widths. If this jet vanishes, the expansion still bounds the
surcharge. Thus even at an exact face the specified cutoff generally
\emph{costs} energy. The available saving in a cap comparison comes
from the discarded side. The full computation, including the profile
identity and the approximate case, is in~\cite{Baker2026CollarDeformation}.

For an approximately matched face take the actual Gaussian chart with
uniform ellipticity relative to \(dt^2+\gamma\). Let
\(\alpha=\|h_0-\gamma\|_{C^2_\gamma}\),
\(\beta=\|h'_0\|_{C^2_\gamma}\) and
\(M_j=\sup\|h''_t\|_{C^2_\gamma}\) on the respective strip.
For \(w_j\le1\) the two normal differentiations produce terms of orders
\(\alpha/w_j^2,\ \beta/w_j,\ M_j\). Taking the square root of the strip
volume explains the powers in
\[
 Z_j=\frac{\alpha}{w_j^{3/2}}+\frac{\beta}{\sqrt{w_j}}+\sqrt{w_j}M_j .
\]
The chiral curvature \(L^2\) norm of each deformed strip is bounded by
\(C\sqrt{\operatorname{vol}_\gamma(S^3)}
 (Z_j+w_j^{3/2}Z_j^2)\). The quadratic term accounts for products of
first derivatives~\cite{Baker2026CollarDeformation}. Shrinking a width
with a fixed nonzero face mismatch can make the estimate worse.
Small raw curvature mass on the strip gives no bound on these marked
jets. A collar selected by another method, for example the physical
annular comparison in~\cite{Baker2026MarkedLimits,Baker2026ChiralMaster},
can be used if its actual metric, markings and estimates supply the
required bounds.

\subsection{From a joined path to a metric on a capped side}

Let \(g_s\) be a smooth path from \(g\) on the joined manifold to a
metric with a round product germ in the specified cut. Define
\[
 H_*=\max_{0\le s\le1}E_+(g_s;M),\qquad
 \Delta_j^\sigma(s)=E_\sigma(g_s;P_j)-m_j^\sigma .
\]
They satisfy \(E_\sigma(g_s;M)=e+\sum_j\Delta_j^\sigma(s)\).
The maximum is taken on the uncut manifold, including both connectors.

\begin{proposition}[Positivity and smooth reassembly]\label{prop:cap}
Assume \(Y([g])>0\) and \(H_*<2T\). The two capped sides are smooth
homotopy four-spheres with positive classes of energies
\(m_j^\sigma+\Delta_j^\sigma(1)\) in sign \(\sigma\). Their marked
connected sum is the original \(M\). If the two positive cap energies
have sum at most \(2T\), at most one capped side can be nonstandard.
If such a side exists it is oriented-diffeomorphic to \(M\).
\end{proposition}
\begin{proof}
On a scalar-flat Yamabe representative, Gauss--Bonnet and signature
give \(E_+=2T+\tfrac14\int|\Ric|^2\,dV\). The first conformal
Laplacian eigenvalue is continuous along a smooth metric path and
its sign equals the Yamabe sign. The strict peak rules out a zero
eigenvalue and therefore keeps all classes positive.
On their unit-volume Yamabe representatives Gauss--Bonnet gives
\(Y([g_s])^2\ge48(2T-H_*)\)
\cite{LeeParker1987Yamabe,Baker2026ChiralMaster}. Adding any finite
length to the round cylinder is a path on the joined manifold at
constant energy. Every compactly supported Yamabe test on a half with
an infinite cylindrical end extends by zero to a sufficiently long
joined manifold and obeys the same positive lower bound. Multiplying
the end metric by \(e^{-2t}\) for large outward arclength \(t\) gives
the radial ball metric in \(r=e^{-t}\). Tests away from the added
center are dense in \(H^1\) in dimension four. This proves positivity
on each smooth cap. Conformal changes preserve both densities and
the added cylindrical and radial portions have zero Weyl curvature.

The middle sphere separates by \(H_1(M)=0\). Van Kampen and
Mayer--Vietoris give trivial fundamental group and zero second
homology on each cap. Duality and the homological Whitehead theorem
then identify each as a homotopy sphere. The finite recognition
criterion \(Y>0,\ E_+\le T\) identifies any such cap with the
standard sphere, including at equality
\cite[Theorems A and B]{ChangGurskyYang2003}%
\cite{Baker2026ChiralMaster}. If the sum is at most \(2T\), at most
one can exceed \(T\). In a standard capped piece its particular cap
ball can be moved to a coordinate ball by Palais's disk theorem.
Its complement is a standard ball. The discrepancy between the
actual boundary identification and a cap parametrization extends
over a standard ball by the \(S^3\) diffeomorphism theorem, using a
collar isotopy and a linear orthogonal map near the center
\cite{Palais1960Extending,Hatcher1983Smale}. Removing the actual caps
and applying the original identification thus gives \(M\) on the
remaining side. This proof also works with the inclusive sum bound
when positivity of both capped metrics is supplied independently.
\end{proof}

There is a useful finite test when neither side is known to be a
disk. Suppose a hypothetical nonstandard \(M\) has initial energy
\(e\) in a positive class and the two positive cap energies are at
most \(m_j^\sigma+\gamma_j^\sigma\), with
\(\gamma_j^\sigma\ge0\) in a fixed sign. If their actual sum is at
most \(2T\), denote the unique nonstandard side by \(s\). Reassembly
and the infimum on \(M\) give
\[
 I^+(M)\le E_\sigma(\widehat g_s)\le m_s^\sigma+\gamma_s^\sigma,
 \qquad
 \min_j m_j^\sigma\le e-I^+(M)+\max_j\gamma_j^\sigma .
\]
Thus a cut can test minimization without deciding in advance which
side would survive. One must price the cap on each possible survivor.
The inclusive test of the actual cap sum applies after their Yamabe
signs are known.

\begin{theorem}[What an exact face forces on each side]\label{thm:mass}
Let \(g\) be a positive Yamabe class representative on \(M\) with
\(e<2T\), and suppose a separating embedded \(S^3\) is conformally
round and umbilic. Let \(\widehat P_j\) be the two sides capped in
their actual markings. Then
\[
 m_j^+=m_j^-=:m_j,\qquad I^+(\widehat P_j)\le m_j\quad(j=L,R).
\]
If \(M\) is nonstandard, precisely one capped side has that smooth
type. On the other side \(m_j\le e-I^+(M)\).
\end{theorem}
\begin{proof}
Fix the smooth source and let the two available subwidths tend to
zero there. If the angular map used to normalize the face differs
from a prescribed cap marking, their transition extends over that
ball as in Proposition~\ref{prop:cap}. This fixes the capped smooth
types throughout the comparison. The cutoff has joined peak at most
\(e+C_g(w_L+w_R)<2T\) for all sufficiently small widths. On the
fixed capped type Proposition~\ref{prop:cap} produces genuine positive
metrics of energy \(m_j^\sigma+O_g(w_j)\). The two chiral integrals
on each closed cap coincide by signature. Taking the limit gives
the mass identity and taking the infimum gives the claimed upper
bound. The two raw masses sum to \(e<2T\), so one is strictly less
than \(T\). At some positive widths that cap is already recognized.
If \(M\) is nonstandard, the other cap is diffeomorphic to \(M\)
by the actual reassembly. Its raw mass is at least \(I^+(M)\), which
gives the last inequality.
\end{proof}
This is a statement about infima on specified smooth types. No
minimizer on either closure is required. The cap energies at positive
width generally differ from the limiting raw masses.

\section{The charged disk as a benchmark}\label{sec:benchmark}

Suppose \(M\) is nonstandard and \(I=I^+(M)<2T\). Start with any
minimizing sequence of smooth unit-volume positive Yamabe metrics
on this fixed smooth type. Write \(e_i=E_\pm(g_i;M)\to I\) and
\(d_i=e_i-I\ge0\). The charged-disk theorem in the master manuscript
supplies, on a further subsequence, smoothly parametrized standard
disks \(D_i\) with a common selected sign \(\sigma_0\) and fixed
\(a,b>0\) such that
\[
 E_{\sigma_0}(g_i;D_i)\ge a,\qquad
 E_{\sigma_0}(g_i;M\setminus\operatorname{int}D_i)\ge b .
\]
These are masses in the original metric on the two whole regions,
including their collar halves~\cite{Baker2026ChiralMaster}.
The theorem gives a disk and a marking. It leaves its jets to be
measured on the supplying source.

Put \(V_i=M\setminus\operatorname{int}D_i\) and
\(m_i^\sigma=E_\sigma(g_i;D_i)\). If a subcritical joined path gives
a round product in this disk marking, the capped exterior is on the
original smooth type: the angular transition between the new cap
and the original parametrized disk extends across a ball. Thus
\begin{equation}\label{eq:finite}
 I\ \le\ E_\sigma(\widehat g_{V,i};M)
     = e_i-m_i^\sigma+\Delta_{V,i}^\sigma(1).
\end{equation}
Here \(\Delta_{V,i}^\sigma(1)\) includes every metric change on the
retained exterior. The complementary change contributes to the
joined peak. A general replacement with a curved cap must also
charge that cap's entire energy. A single strictly positive-width
construction with net saving \(m_i^\sigma-\Delta_{V,i}^\sigma(1)>d_i\)
would violate~\eqref{eq:finite} and exclude the hypothesized type.
Equality is consistent with attainment of \(I\). Independently, an
attained positive energy at most \(T\) recognizes the type, and the
master proves recognition for the inclusive infimum bound
\(I\le T+\varepsilon_*\) with a universal, non-effective
\(\varepsilon_*>0\)~\cite{Baker2026ChiralMaster}.

For an exact conformally round umbilic boundary, apply
Theorem~\ref{thm:mass} to the cut and identify its capped exterior
with \(M\). This gives
\(m_i^+=m_i^-\le d_i\) at this fixed finite index with \(e_i<2T\).
Since \(m_i^{\sigma_0}\ge a\) and \(d_i\to0\), every sufficiently
late selected disk fails the exact-face hypothesis. Widths are taken
to zero at each individual index before making this conclusion. There
is also a conditional approximate exclusion. Suppose the physically
marked bicollars, in smooth conformal gauges on a fixed slab, are
smoothly precompact with nondegenerate limits. If the integral of
\(|\operatorname{Cotton}(h_i)|+|K_i^\circ|^3\) on the central
spheres tends to zero, the limit is conformally round and umbilic.
The limit chart and subwidths with bounded ratio of \(C^2\) metric
error to squared width then produce \(o(1)\) complete-side and joined
costs, with clearance for any separately protected regions. Equation
\eqref{eq:finite} gives a contradiction at a finite late index.
Thus the defect has positive lower limit on any such precompact
charged subsequence~\cite{Baker2026CollarDeformation}. Disk selection
alone gives neither compactness nor a numerical defect bound.

\paragraph{Why changing only a thin strip can be insufficient.}
Keep the outer disk \(Q_i\) enclosing the original bicollar inside
a larger parametrized disk \(B_i\), and put \(q_i^\sigma=E_\sigma(g_i;Q_i)\).
The master's simultaneous measure division at fixed depth \(k\ge1\)
gives raw mass
\[
 b_i^\sigma=E_\sigma(g_i;B_i)
  =(1-2^{-k})e_i+2^{-k}q_i^\sigma-r_i^\sigma,\qquad r_i^\sigma\to0 .
\]
This follows by halving the remaining exterior mass at each step
up to a chosen tolerance, with all \(k\) errors included in
\(r_i^\sigma\). For a matched endpoint \(h_i\) with joined peak
\(H_i<2T\), set \(S_i^\sigma=b_i^\sigma-E_\sigma(h_i;B_i)\).
The new capped exterior still has energy at least \(I\) in each
sign. Hence
\begin{align}
 S_i^\sigma &\ge b_i^\sigma+I-H_i>b_i^\sigma+I-2T,\label{eq:saving}\\
 \Delta_{V_i^B}^\sigma&\ge b_i^\sigma-d_i,\qquad
 \Delta_{V_i^B}^\sigma\ge S_i^\sigma-d_i,\qquad
 V_i^B=M\setminus\operatorname{int}B_i . \nonumber
\end{align}
Indeed \(E_\sigma(h_i;B_i)\le H_i-I\), and subtracting the two
endpoint side identities gives
\(\Delta_{V_i^B}^\sigma=S_i^\sigma+E_\sigma(h_i;M)-e_i\),
where the positive joined endpoint has energy at least \(I\).
If \(Q_i\) is fixed up to conformal change, then
\(H_i\ge I+\max_\sigma q_i^\sigma\).
Since \(I>T+\varepsilon_*\) on the hypothetical type, \(k\) can be
fixed large enough that \((2-2^{-k})I-2T>0\). The right side of
\eqref{eq:saving} then has a positive lower limit in both signs.
A change confined to a disk-side strip of vanishing raw mass cannot
pay it. The estimate gives the saving required of an actual path.
Constructing such a path remains a variational problem
\cite{Baker2026ChiralMaster,Baker2026FixedDescent,Baker2026HandleCuts}.

\paragraph{A curved filling and the source of positivity.}
A second comparison uses the actual exterior of the charged disk
without stretching a joined neck. In the positive Yamabe metric let
\(h,K,H\) be the boundary metric, outward disk shape and its trace.
With \(G=R_h\mathrm{Id}/2-\Ric_h^\sharp>0\), set
\(S=\sqrt{\det G}\,G^{-1}\). If \(S\) is Codazzi, it is realized
by a convex Euclidean ball with that boundary metric. If also
\(S^\circ=K^\circ\), a cap conformal factor matches the two full
first jets. Smoothing on the cap side at fixed source then costs
\(O_i(w)\) in each sign and preserves the entire exterior.
When \(H\le\mathrm{tr}\,S\) the joined endpoint form is positive
by the exterior scalar curvature, flat Dirichlet form and nonnegative
boundary contribution. Thus \(m_i^+=m_i^-\le d_i\)
\cite{Baker2026CollarDeformation}. This excludes these simultaneous
boundary measurements on every sufficiently late original charged
disk, without assuming compactness of its marked collars. A natural
candidate from an immersed height illustrates the difficulty:
on a transverse nonround section of positive definite outward shape
the flat section's trace-free shape differs from the induced
four-dimensional one wherever it is nonzero.
An independently positive deformation supported in the disk interior
can save at most \(d_i\). If its support is disjoint from the complete
seam change, the capped exterior has exactly the same energy as for
the seam change alone \cite{Baker2026CollarDeformation}.
The selected minimizing sequence still lacks a measured fit or a
saving that pays the complete exterior cost.

\section{Selection, termination and outstanding challenges}

Choose a presentation with the least middle count \(n\) on a
hypothetical nonstandard \(M\). The count exists among finite Morse
presentations, and Section~\ref{sec:handles} gives \(n\ge2\). A sufficient
selection theorem would give at least one of the following conclusions
for every such remaining type and presentation.
\begin{enumerate}
\item A relative smooth move yields a presentation of \(M\) with
strictly smaller \(n\), counted after all temporary stabilizations,
or identifies the actual closed manifold with a recognized one.
\item A geometric test recognizes \(M\) directly. Examples are a
dominated recognizing prefix in Section~\ref{sec:handles} and an
actual positive class of energy at most \(T\).
\item The type has positive subcritical entry, and a physical repair
of the supplied regions at one finite index gives a positive
competitor on this \(M\) of energy strictly below \(I^+(M)\).
\end{enumerate}
The first outcome contradicts minimality or identifies the type.
The second identifies it independently of the count. The third
contradicts the definition of the infimum. This conditional implication
explains why a net cancellation or a single strict comparison would
suffice. A proof of the selection hypothesis would be a further
theorem.

\paragraph{The first unresolved handle presentations.}
For a winding-one knot with three geometric belt crossings, a
suitable embedded bigon permits a cancellation with movable residual
tubes. If that disk is missing, one can seek a different framed slide,
the dual belt-circle isotopy, or domination on the opposite cut.
The first uncovered mixed class has none of these certificates in
the tested descriptions. In the one-sided class with two components,
or after the three-to-two reduction above, the residual two-component
link can fail to have any identified recognizing component or trace.
With several lower 1-handles the penultimate test still applies
when its sublevel is simply connected. In the other cases a different
certified ordering or relative move is needed. The unknot and band arguments do not select one on an arbitrary
diagram. A positive outcome for one particular presentation would
advance exactly the class on which that outcome can be repeated.

\paragraph{The first unresolved metric comparison.}
On a late charged disk an exact round umbilic face would already
contradict minimization. A quantitative replacement must therefore
control the actual nonzero defect or the degeneration of its marked
chart, or modify a larger region while paying the full two-sided
cost. One target is a nonconformal repair that accounts for the protected
core and gives a positive exterior cap with energy strictly below \(I\)
at one finite index, contradicting~\eqref{eq:finite}. Another is a
different spherical cut with enough raw mass on
each side to violate the finite cap inequality. Relating a disk to
a handle level requires an identification and metric estimates on
the same source.

\paragraph{Entry on an arbitrary smooth type.}
Every \(M\) supplies a finite handle presentation. Universal selection
of a reducing move is open. For an independent geometric entry, a
Euclidean hypersurface immersion of induced \(E_+<8\pi^2\) on a
nonstandard type would give more than three quarters of directions
with exactly two middle handles. One recognized marked height would
then suffice~\cite{Baker2026ImmersionEnergy,Baker2026MarkedHeights}.
Finding such an immersion on every remaining type, finding a recognized
height even when it exists, and producing an intrinsic positive
class of energy below \(2T\) when following the charged-disk route are
distinct outstanding problems. The induction and the metric comparison
give precise conclusions once one of their geometric inputs is found.

\begingroup
\raggedright
\bibliographystyle{unsrturl}
\bibliography{references}
\endgroup
\end{document}

